\section{Producing finite algebra certificates by bounded linear algebra}
\label{sec:macaulay}

The inspected packet \texttt{ipra\_presentations\_to\_witnesses} already proves
that finite multiplication identities yield a complete list of rational
zeros~\cite{repo-witnesses}. Its constructors start with separated annihilator identities or a
nonzero plane norm. We add a more general producer: a single finite linear
algebra test at supplied degree caps constructs all required identities.
This is a new interface relative to the inspected packets, not a claim of
historical novelty for polynomial linear algebra. It is related to
Macaulay and border-basis methods \cite{mourrain2008}. We do not need a quotient basis, a
normal form, or commuting matrices; those are stronger conclusions than
the finite-image argument needs.

\subsection{The input and the exact acceptance test}

Fix $r\geq1$ and a finite list $f_1,\ldots,f_s$ in
$\mathbb Q[X_1,\ldots,X_r]$. A polynomial is a finite coefficient table;
its entries are reduced signed numerators and positive denominators.
Let $d,D$ be supplied natural numbers with $D\geq d+1$.
Write $P_e$ for the rational vector space with monomial list
\[
 B_e=\{X^\alpha: \alpha\in\mathbb N^r,\ |\alpha|\leq e\},
 \qquad N_e=|B_e|.
\]
Here $P_e$ is only convenient notation for coefficient vectors of length
$N_e$; it does not introduce a new object-theory sort.
Enumerate $B_e$ by scanning the $(e+1)^r$ exponent tuples with coordinates
at most $e$ and retaining those of total degree at most $e$.

Inside $P_D$, form the finite list
\[
 S_D=\{X^\alpha f_i:|\alpha|+\deg f_i\leq D\}.
\]
Zero input polynomials may be omitted. No generator of degree above $D$
contributes a column. Let $A$ be the rational matrix whose columns are,
first, $B_d$ and, second, the members of $S_D$, padded by zero to the
coefficient coordinates $B_D$. For each monomial $b$ of total degree
$d+1$, solve the finite linear system
\begin{equation}
                       A c_b=[b]_{B_D}.
\label{eq:border-linear-system}
\end{equation}
Accept precisely when every one of these systems is solvable.
The required solutions are produced by rational Gaussian elimination.
There is no enumeration of arbitrary rational coefficients and no
coefficient-height guessing.

\begin{lemma}[A bounded solver with success and failure evidence]
\label{lem:bounded-linear-solver}
Given $A\in\mathbb Q^{m\times n}$ and $v\in\mathbb Q^m$, a fixed
primitive recursive program returns either $c\in\mathbb Q^n$ with
$Ac=v$, or $u\in\mathbb Q^m$ with $uA=0$ and $uv\ne0$.
Its finite arithmetic correctness can be proved in iPRA.
\end{lemma}

\begin{proof}
Perform Gaussian elimination on the augmented matrix $[A\mid v]$, using
only its first $n$ columns for pivot selection. Scan columns in order; in
each column scan the unused rows for the first nonzero entry. If there is
none, continue. Otherwise swap that row into the next pivot position,
divide by its nonzero pivot, and subtract its suitable multiples from all
other rows. Carry along an initially identity matrix $U$, subject to the
same row operations. The invariant is
\[
             [R\mid w]=U[A\mid v].
\]
There are at most $m$ pivots and exactly $n$ column scans. Every pivot
operation and every coefficient update has an explicitly bounded finite
index set. The invariant and the maintained pivot pattern are bounded
conjunctions of rational equalities and inequalities.

If a final row of $R$ is zero while the corresponding entry of $w$ is
nonzero, the same row of $U$ is the desired $u$. Indeed $uA=0$ and
$uv\ne0$. A putative solution $c$ would give $uv=uAc=0$, so this is
a genuine incompatibility certificate for this linear system.
Otherwise set every free variable to zero and every pivot variable to
the corresponding right-hand entry. The reduced pivot pattern then gives
$Rc=w$. Since row swaps, nonzero row rescalings, and row additions are
invertible, their inverse operations show $Ac=v$. One may track a second
matrix for the inverse operations; both products being identity are finite
matrix equalities. Alternatively replay the finite operation log backwards.

Finite sums, rational arithmetic, canonical normalization, finite matrix
updates, and bounded scans are primitive recursive. Rational inversion is
used only after an exact nonzero test. Gcd normalization may itself be
implemented by a bounded divisor scan, so no unbounded termination claim
is needed. Booleanize the bounded matrix invariants, apply quantifier-free
induction to the fixed loops, and afterwards generalize the free entries
of $A,v,c$. This gives the stated iPRA proof route without induction on a
formula quantifying over all possible solutions.
\end{proof}

Failure of one instance of \eqref{eq:border-linear-system} means that the
supplied caps did not establish closure. It says nothing by itself about
whether the polynomial system has rational zeros, or about whether larger
caps would succeed. The returned separating vector certifies exactly this
finite failure.

\subsection{From closure to a multiplication certificate}

Put $B=B_d=(b_1,\ldots,b_N)$, ordered so that $b_1=1$.
An accepted solution gives, for each degree-$d+1$ monomial $b$, explicit
rational coefficients $c_{bi}$ and explicit polynomials $h_{b\ell}$ with
\begin{equation}
       b-\sum_{i=1}^N c_{bi}b_i
                  =\sum_{\ell=1}^s h_{b\ell} f_\ell .
\label{eq:border-identity}
\end{equation}
The $h_{b\ell}$ are assembled from the coefficients of the columns
$X^\alpha f_\ell$ in $c_b$. Every such summand has degree at most $D$.
Equation \eqref{eq:border-identity} is therefore checked coefficient by
coefficient in the same finite coordinate space $P_D$.

For each coordinate $k$ define an $N\times N$ matrix $M_k$ columnwise.
If $X_kb_j$ has degree at most $d$, use its exact unit coordinate vector
in $B$. Otherwise use the coefficient vector furnished by
\eqref{eq:border-identity}. This produces identities
\begin{equation}
        X_k b_j-\sum_i(M_k)_{ij}b_i
                       =\sum_\ell h_{kj\ell}f_\ell.
\label{eq:border-action}
\end{equation}
The unit identity is simply $b_1=1$.

\begin{theorem}[A complete producer at supplied degree caps]
\label{thm:border-producer}
There is one primitive recursive program, uniform in the numbers of
variables and equations, their finite coefficient tables, and $d,D$,
which decides the finite closure test and, upon acceptance, produces
the explicit identities \eqref{eq:border-action}. These identities yield
a duplicate-free finite list $L$ of exactly the rational common zeros:
\[
 p\in L\quad\Longleftrightarrow\quad
 f_1(p)=\cdots=f_s(p)=0\qquad(p\in\mathbb Q^r).
\]
The soundness and completeness implication is an iPRA theorem after
the finite arithmetic contracts in this proof are compiled.
No independence or commutation assertion about the $M_k$ is assumed.
\end{theorem}

\paragraph{The finite divisibility lemma used in the proof.}
We record the elementary arithmetic behind the coprimality steps, rather
than treating them as unanalysed number-theory calls. Given $a,b\in\mathbb N$
with $b>0$, initialize remainders $(r_0,r_1)=(a,b)$ and coefficient rows
$(s_0,t_0)=(1,0)$, $(s_1,t_1)=(0,1)$. At an active step, put
\[
 q=r_0\mathbin{\mathrm{div}}r_1,\qquad
 r_2=r_0-q r_1,\qquad
 (s_2,t_2)=(s_0-q s_1,t_0-q t_1),
\]
and replace rows zero and one by rows one and two. Freeze once $r_1=0$.
The positive second remainder strictly decreases at every active step,
so exactly $b$ fuel steps suffice. The invariants
$r_i=s_i a+t_i b$ and preservation of the common divisors of the two
remainders are finite arithmetic statements. The latter follows in both
directions from $r_2=r_0-q r_1$. A possible common divisor is kept as a
free parameter during induction and generalized afterwards; alternatively
one scans divisors up to $\max(a,b)$. Thus no unbounded universal
quantifier enters the induction formula. The final positive remainder is
$g=\gcd(a,b)$ and its row supplies $s a+t b=g$.
When one input is zero, the same computation or the immediate zero case
supplies the corresponding identity. Signs are incorporated by changing
the coefficient of a negative input.

Consequently $\gcd(|a|,b)=1$ supplies $s a+t b=1$. If $b\mid a w$,
multiplication by $w$ proves $b\mid w$. Folding the same bounded
extended-Euclid construction along a nonzero integer vector supplies
coefficients $e_i$ with $\sum_i e_i u_i=g$, where $g$ is the gcd of its
absolute coordinates. On dividing the vector by $g>0$, this becomes
$\sum_i e_i(u_i/g)=1$. These are fixed PR computations with bounded
loop invariants, even though their integer coefficients can be large.
They justify both coprimality cancellation and the primitive-vector
Bezout combination below.

\begin{proof}
Construction of the identities was just given. It remains to prove that
they supply a finite and complete coordinate grid. This repeats the
finite-image argument from the inspected packet so that the present
result is self-contained.

For a common zero $p$, evaluation of \eqref{eq:border-action} gives
\[
       vM_k=p_kv,\qquad v=(b_1(p),\ldots,b_N(p)),\qquad v_1=1.
\]
Thus $v$ is nonzero. Multiply $v$ by the product of its denominators,
then divide the resulting integer vector by the gcd of its coordinates.
This gives a primitive nonzero integer row vector $w$ satisfying the
same eigenvector equations.

Choose a positive common denominator $\Delta_k$ for $M_k$ and put
$A_k=\Delta_kM_k\in\mathbb Z^{N\times N}$. We have
$wA_k=\lambda_kw$ with $\lambda_k=\Delta_kp_k$ rational.
Write $\lambda_k=a/b$ in lowest terms, $b>0$. The coordinate equations
give $b\mid a w_j$, hence $b\mid w_j$ by coprimality of $a,b$.
A Bezout combination of the primitive coordinates of $w$ equals one,
so $b\mid1$ and $b=1$. Consequently $\lambda_k$ is an integer.

Let
\[
        R_k=\max_j\sum_i|(A_k)_{ij}|.
\]
Choose a first coordinate with $|w_j|=W=\max_i|w_i|>0$. Then
\[
 |\lambda_k|W
 =\left|\sum_iw_i(A_k)_{ij}\right|
 \le W\sum_i|(A_k)_{ij}|\le WR_k.
\]
Cancellation of the positive $W$ yields $|\lambda_k|\le R_k$.
Every common zero therefore lies in the explicit grid
\begin{equation}
       \prod_{k=1}^r
          \{z/\Delta_k:z\in\mathbb Z,\ -R_k\le z\le R_k\}.
\label{eq:border-grid}
\end{equation}
Each coordinate in reduced form has denominator dividing $\Delta_k$
and numerator of absolute value at most $R_k$.

Enumerate the $\prod_k(2R_k+1)$ grid tuples by mixed-radix decoding,
evaluate every $f_\ell$ exactly, and retain precisely the common zeros.
The retention test proves soundness and the grid argument proves
completeness. Fixed nonzero denominators make distinct grid numerators
produce distinct rational coordinates, so the list has no duplicates.

This argument used a separate eigenvector equation for each matrix.
The matrices need not define an algebra structure on the whole ambient
vector space: redundant spanning elements may permit noncommutation on
directions which no evaluation vector detects. The ordinary theory of
border \emph{bases} imposes further conditions to obtain a unique normal
form and a true basis; the present certificate does not assert these
stronger conclusions.

For the uniform PR claim, $N_e\le(e+1)^r$, and the number of generator
columns is at most $s(D+1)^r$. Coefficient generation, each pivot loop,
all identity checks, denominator products, gcd folds, and the final grid
scan have finite bounds computed from input data. Arbitrarily large
intermediate coefficients do not obstruct primitive recursion: each
finite update is a fixed PR operation and the iteration counts are fixed
PR terms. Uniform polynomial evaluation is a fixed bounded interpreter
for coefficient tables and powers; it is not a universal evaluator for
arbitrary primitive recursive program codes.

For the object-theory proof, keep $p$ and its coded evaluation vector as
free parameters. Use bounded conjunctions for coefficient and matrix
identities, convert them to Boolean equalities, and induct on the finite
elimination, normalization, maximum, and grid loops. The correctness
theorem of the final filter is proved for this arbitrary $p$ and then
universally generalized. No global assertion that every accepted formal
proof is true is used.
\end{proof}

\begin{lemma}[Finite span without a canonical reduction]
\label{lem:action-finite-span}
Suppose $b_1=1$ and the identities \eqref{eq:border-action} hold for a
fixed finite family $b_1,\ldots,b_N$. Every polynomial is congruent,
modulo the ideal generated by the $f_\ell$, to a rational linear
combination of this family. An explicit PR procedure produces both the
coefficient vector and the ideal-membership identity. It need not
produce a unique vector.
\end{lemma}

\begin{proof}
First consider a monomial $X^\alpha$. Choose the fixed ordered word
consisting of $\alpha_1$ copies of $X_1$, followed by $\alpha_2$ copies
of $X_2$, and so on. For the empty word, use coefficient vector $c=e_1$
and witnesses $h_\ell=0$, so that
\[
                  1-\sum_j c_jb_j=\sum_\ell h_\ell f_\ell.
\]
Maintain, for the word prefix whose monomial is $m$, the explicit identity
\[
                  m-\sum_j c_jb_j=\sum_\ell h_\ell f_\ell.
\]
On appending $X_k$, put
\[
 c'_i=\sum_j(M_k)_{ij}c_j,\qquad
 h'_\ell=X_kh_\ell+\sum_j c_jh_{kj\ell}.
\]
Multiply the previous identity by $X_k$ and substitute
\eqref{eq:border-action}. The result is exactly the maintained identity
with $m'=X_km$, $c'$, and $h'$. After $|\alpha|$ steps it gives the
required reduction of the monomial. Reduce each of a polynomial's
finitely many terms in this fixed order, multiply its coefficient vector
and witnesses by the term coefficient, and add them.

The input coefficient tables bound every loop; each update is a finite
polynomial or matrix operation. Coefficientwise identities are the
quantifier-free induction invariants. Different variable orders may give
different coefficient vectors, since the family may be dependent. The
fixed order suffices for existence of a reduction, so commutation of the
matrices on the full coordinate space is unnecessary. In ordinary
quotient notation, the classes of the $b_j$ therefore span
$\mathbb Q[X_1,\ldots,X_r]/(f_1,\ldots,f_s)$ and its dimension is at
most $N$.
\end{proof}

\subsection{Two worked examples and a rejected cap}

Consider
\[
 f_1=XY-1,\qquad f_2=X^2-Y.
\]
For $d=1,D=2$, the list $B=(1,X,Y)$ and the generator columns cannot
span $Y^2$: its coefficient is zero in every column. The coordinate
functional extracting the $Y^2$ coefficient is a failure certificate.
Increasing $D$ to three succeeds, because
\[
       X^2=Y+f_2,\qquad XY=1+f_1,
       \qquad Y^2=X+Xf_1-Yf_2.
\]
The matrices in the order $(1,X,Y)$ are
\[
 M_X=\begin{pmatrix}0&0&1\\1&0&0\\0&1&0\end{pmatrix},
 \qquad
 M_Y=\begin{pmatrix}0&1&0\\0&0&1\\1&0&0\end{pmatrix}.
\]
Here $\Delta_X=\Delta_Y=R_X=R_Y=1$.
Exactly nine grid points are tested; the complete rational solution list
is $\{(1,1)\}$. Independently, the equations imply $X^3=1$ and hence
$(X-1)(X^2+X+1)=0$. Since
$X^2+X+1=(X+1/2)^2+3/4>0$ over the rationals, $X=Y=1$.
The failed cap did not mean that the system was inconsistent.

For a denominator-sensitive example let
\begin{align*}
 f_1&=2(X+Y)^2-(X+Y),\\
 f_2&=3(X-Y)^2-(X-Y),\\
 f_3&=X^2-Y^2.
\end{align*}
The caps $d=1,D=2$ suffice. Exact identities are
\begin{align*}
 X^2-\frac{5X+Y}{24}&=\frac18f_1+\frac1{12}f_2+\frac12f_3,\\
 XY-\frac{X+5Y}{24}&=\frac18f_1-\frac1{12}f_2,\\
 Y^2-\frac{5X+Y}{24}&=\frac18f_1+\frac1{12}f_2-\frac12f_3.
\end{align*}
They produce
\[
 M_X=\begin{pmatrix}0&0&0\\1&5/24&1/24\\0&1/24&5/24\end{pmatrix},
 \quad
 M_Y=\begin{pmatrix}0&0&0\\0&1/24&5/24\\1&5/24&1/24\end{pmatrix}.
\]
Taking both denominators to be $24$ gives both column bounds $24$.
The grid has $49^2=2401$ points. The complete list is
\[
              (0,0),\qquad (1/4,1/4),\qquad(1/6,-1/6).
\]
Indeed put $u=X+Y,v=X-Y$. Then $u\in\{0,1/2\}$,
$v\in\{0,1/3\}$, and $uv=0$. These three possibilities give exactly
the displayed points. The exact reference program independently produces
the reduction identities and filters the full grid.

The program also accepts $X^2-1=0$, $Y-X=0$ using $B=(1,X,Y)$ and
$d=1,D=2$. This spanning family is redundant and its produced matrices
do not commute, yet the complete list is correctly
$\{(-1,-1),(1,1)\}$. For $X^2-2=0$, $Y-X=0$ it accepts the same
caps and returns an empty rational list. These examples check the precise
difference between an algebraic certificate, independence of its
generators, and existence of rational points.

\subsection{What has been resolved and what has not}

The new step removes the need to guess coefficient heights or hand-build
the action identities once finite degree caps have been supplied.
It also provides independently checkable negative evidence when those
caps fail. It covers finite polynomial systems in any input dimension,
with equality decided on exact rational coefficient tables.

No unbounded loop increasing $d,D$ is included in the PR theorem.
General termination or a uniform cap would require an additional
effective algebraic argument. Even success for every zero-dimensional
ideal would not prove general Faltings: a positive-dimensional curve
can have finitely many rational points while its coordinate algebra is
infinite-dimensional. This is the precise boundary relevant to the
inspected Siegel and Faltings routes~\cite{repo-siegel,repo-faltings}.

For completeness, take $F=Y^2+X^2+1$. There are no rational zeros, since
$X^2+Y^2\geq0$. Nevertheless monic division in $Y$ gives each polynomial
a remainder $a(X)+b(X)Y$ modulo $F$. The division algorithm subtracts
successive multiples of the monic leading term and decreases $Y$-degree
at each active step, so an explicit degree fuel makes it primitive
recursive. The remainder is unique: if a nonzero $H(X,Y)$ existed with
$HF$ of $Y$-degree below two, its leading $Y$-coefficient would remain
the nonzero leading coefficient of $H$, and
$\deg_Y(HF)=\deg_Y H+2\geq2$, a contradiction. Thus the quotient is
the free $\mathbb Q[X]$-module on $1,Y$. In particular its elements
$1,X,X^2,\ldots$ are linearly independent over $\mathbb Q$.

A finite action certificate of size $N$ would contradict this using only
finite data. Lemma~\ref{lem:action-finite-span} would reduce
$1,X,\ldots,X^N$ to $N+1$ vectors of length $N$. Rational Gaussian
elimination supplies a nontrivial rational dependence among those
vectors. Combining their explicit ideal-membership witnesses yields a
nonzero $q(X)\in(F)$. Its remainder under the preceding monic division
is $q(X)$ itself, whereas every member of $(F)$ has zero remainder.
This contradiction shows that a full coordinate multiplication
certificate cannot exist for that sole equation, despite its empty
rational zero set. A finite algebra section or a supplied finite cover
by such sections is a valid use of the producer; existence of the cover
remains a separate theorem.

\paragraph{Lean implementation targets.}
Use finite coefficient lists over a coded rational type. Define
\texttt{monomials}, \texttt{macaulayColumns}, \texttt{rrefWithTrace},
\texttt{closureCertificate}, \texttt{checkAction}, and
\texttt{rationalGrid} as fuelled list programs. The main successive lemmas
are the row-operation invariant, the solution-or-separator specification,
coefficientwise construction of \eqref{eq:border-action}, the primitive
eigenvector denominator lemma, the column bound, and filter completeness.
The host Lean statements are ordinary arithmetic lemmas; the object
targets must explicitly conclude \texttt{Prf PRAAx} for their encoded
Boolean contracts. The supplied Python checks establish only exact
finite examples. They are not a Lean proof or an iPRA proof checker.
